\section{The endpoint inequality at norm 89}
\label{sec:endpoint}

\begin{theorem}[The norm-$89$ endpoint inequality]
\label{thm:endpoint-89}
Let $q=\varepsilon\chi(w)2$ be any word in \eqref{eq:regular-source}, put
$v=M(q)e_1$, and assume $\ord_{89}(\|v\|^2)=1$. Then its complete norm-$89$
flip satisfies $\Delta\geqslant0$.
\end{theorem}

The proof replaces the infinitely many admissible words by the finitely
many states of a reduction graph, and then exhibits on that graph an
integer potential making every step of the cost comparison nonnegative.
The finiteness of the graph is not an empirical stopping observation: it
is forced by the determinant.

\subsection{From row subtraction to weighted transitions}

Suppose a nonnegative integral matrix $S$ has two incomparable rows:
neither dominates the other componentwise. After multiplying by $M(j)^2$,
subtract a dominated row from the dominating row until the rows are again
incomparable. These are the elementary row operations
$$
 (a,b;c,d)\longmapsto(a-c,b-d;c,d)
 \quad\hbox{or}\quad
 (a,b;c,d)\longmapsto(a,b;c-a,d-b).
$$
The notation lists the two rows separated by a semicolon. Every such
subtraction is forced for \emph{every} positive continuation vector:
componentwise domination implies domination after multiplication by that
vector. The rows remain nonnegative, and the sum of the entries decreases,
so this reduction terminates. If $k$ unit subtractions are emitted, assign
the transition the weight
\begin{equation}
 w(S,j)=k-2j.
 \label{eq:edge-weight}
\end{equation}
The subtracted $2j$ is exactly the source cost of the paired digit $jj$.

Row exchange has no effect on Euclidean cost. We therefore identify a
matrix with its row exchange and choose the representative with positive
determinant. Every reduced representative takes the form
\begin{equation}
 S=\begin{pmatrix}r+u&b\\r&b+v\end{pmatrix},\qquad
 r,b\geqslant0,\quad u,v\geqslant1.
 \label{eq:crossed-state}
\end{equation}

\begin{lemma}[State invariants]
\label{lem:state-invariants}
Every state reachable from a seed has the crossed form
\eqref{eq:crossed-state}, content one, and determinant
\begin{equation}
 rv+bu+uv=89^2.
 \label{eq:finite-state-bound}
\end{equation}
In particular $r,b,u,v\leqslant89^2$, so only finitely many states occur.
\end{lemma}

\begin{proof}
A seed arises from $R_\nu M(\varepsilon)$ by reading paired letters and
performing row operations. The digit matrices are unimodular and
$\det R_\nu=-89^2$, so a seed has determinant $\pm89^2$; the choice of the
representative with positive determinant fixes the sign, and expanding
\eqref{eq:crossed-state} gives \eqref{eq:finite-state-bound}. Row
subtraction, row exchange and sign change are unimodular, and so are the
digit matrices acting on the right; each therefore preserves both the
determinant up to sign and the greatest common divisor of the entries.
Since a seed has content one, so does every state reachable from it.
Finally \eqref{eq:finite-state-bound} writes $89^2$ as a sum of three
nonnegative products, and $u,v\geqslant1$, so each of $r,b,u,v$ is at most
$89^2$.
\end{proof}

If the input ends in state $S$, its remaining vector is $Se$, where
$e=(2,1)^{\mathsf T}$. The terminal charge is
\begin{equation}
 f(S)=\ell(Se)-1.
 \label{eq:terminal-charge}
\end{equation}
We accept the terminal exactly when $\cont{Se}=89$. Unimodular row
operations preserve vector content, so \cref{lem:primitive-flip} makes
this the exact primitive endpoint test. The final one subtracted in
\eqref{eq:terminal-charge} is the source cost of its last digit $2$.

\subsection{Resolving every initial sign}

The first row of each $R_\nu$ in \eqref{eq:reflections} is positive, but
the second has opposite signs. Starting with $R_\nu M(\varepsilon)$,
read paired letters until the second row has one weak sign componentwise.
At that point its sign is fixed for every positive continuation. Change
its sign if necessary, perform the forced row reduction, and enter the
finite graph.

This prefix procedure has eleven leaves. \Cref{tab:seeds} gives every
leaf state and its seed weight. A seed that consumes the paired prefix
$p$ and emits $k_0$ subtractions has weight
\begin{equation}
 w_0=k_0-\varepsilon-2\sum_{j\in p}j.
 \label{eq:seed-weight}
\end{equation}
The optional digit is again interpreted as zero when absent.

\begin{table}[tb]
\centering
\caption{Complete sign-resolution seeds. A state $(a,b;c,d)$ has the
displayed first and second rows and positive determinant $89^2$.
The prefix consists of paired input letters. These eleven leaves cover
every continuing word after at most two such letters.}
\label{tab:seeds}
\begin{tabular}{ccclr}
\toprule
Branch & $\varepsilon$ & Prefix & State $S_0$ & $w_0$\\
\midrule
$R_0$ & $\varnothing$ & $1$ & $(121,41;37,78)$ & $-1$\\
$R_0$ & $\varnothing$ & $2$ & $(223,10;33,37)$ & $0$\\
$R_0$ & $1$ & $\varnothing$ & $(119,39;41,80)$ & $-1$\\
$R_0$ & $2$ & $\varnothing$ & $(158,39;121,80)$ & $-2$\\
$R_1$ & $\varnothing$ & $1$ & $(195,37;2,41)$ & $0$\\
$R_1$ & $\varnothing$ & $21$ & $(68,33;67,149)$ & $10$\\
$R_1$ & $\varnothing$ & $22$ & $(119,1;52,67)$ & $9$\\
$R_1$ & $1$ & $1$ & $(43,2;17,185)$ & $4$\\
$R_1$ & $1$ & $2$ & $(127,43;120,103)$ & $0$\\
$R_1$ & $2$ & $1$ & $(233,20;35,37)$ & $3$\\
$R_1$ & $2$ & $2$ & $(271,23;68,35)$ & $7$\\
\bottomrule
\end{tabular}
\end{table}

The prefix tree has sixteen nodes, so besides the eleven leaves there are
five words that stop earlier. They are the empty paired prefix with absent
optional digit in each branch, the empty prefix with optional digit $1$ or
$2$ in branch $R_1$, and prefix $2$ with absent optional digit in branch
$R_1$. Appending the final digit $2$ gives raw contents
$1,1,1,1,1$, respectively, so none is accepted. The seed table therefore
covers every primitive source in the theorem, including short words, and
no infinite sign-resolution branch is left over.

\subsection{The potential and the proof}

Starting from the eleven seeds, close the state set under the two
transitions. The next statement describes the resulting graph. It is a
finite assertion about integers, and it is the one place in this section
where we appeal to a computation.

\begin{proposition}[The closed graph and its potential]
\label{prop:closed-graph}
The closure of the eleven seeds of \cref{tab:seeds} under the two
transitions has $11420$ states and $22840$ labeled edges, of which
exactly $116$ states satisfy the primitive terminal condition. Let $\lambda(S)$
denote the least weighted length of a path from a seed to $S$, each seed
contributing its own weight $w_0$. Then
\begin{equation}
 \begin{aligned}
 w_0-\lambda(S_0)&\geqslant0 &&\text{at all eleven seeds},\\
 w(S,j)+\lambda(S)-\lambda(S')&\geqslant0 &&\text{on all $22840$ edges},\\
 \lambda(S)+f(S)&\geqslant0 &&\text{at all $116$ primitive terminals}.
 \end{aligned}
 \label{eq:potential-inequalities}
\end{equation}
\end{proposition}

\begin{proof}
By \cref{lem:state-invariants} the state set is finite, and every quantity
appearing in \eqref{eq:potential-inequalities} is an integer; the assertion
is therefore a finite list of integer inequalities. We have verified it by
computer, in exact integer arithmetic: the state set is closed under the
two transitions, the minimum defining $d$ is attained and computed at
every state, and the three families are then checked on the closed graph.
The verification was carried out twice, independently, and the resulting
table of values of $\lambda$ accompanies this article as supplementary
material.
\end{proof}

The values of $\lambda$ may be negative; it is their transformed combinations in
\eqref{eq:potential-inequalities} that are nonnegative.

\begin{proof}[Proof of \cref{thm:endpoint-89}]
Let $q$ be an accepted word, with seed state $S_0$ and terminal state
$S_f$. Every operation emitted along its path is an actual unit
subtraction on the reflected continuation: the two entries of $Sc$ can
never be equal at an emitted step, since the dominating row minus the
dominated one would then annihilate the positive continuation vector $c$,
making a row of $S$ vanish against $c$ and forcing $\det S=0$, against
\eqref{eq:finite-state-bound}; the unconsumed terminal cost is
\eqref{eq:terminal-charge}, and the consumed source charges are exactly
\eqref{eq:seed-weight} and \eqref{eq:edge-weight}. Consequently
\begin{equation}
 \begin{aligned}
 \Delta
 &=w_0+\sum_e w(e)+f(S_f)\\
 &=\bigl(w_0-\lambda(S_0)\bigr)
   +\sum_{e:S\to S'}\bigl(w(e)+\lambda(S)-\lambda(S')\bigr)
   +\lambda(S_f)+f(S_f),
 \end{aligned}
 \label{eq:exact-endpoint-sum}
\end{equation}
the inserted potentials telescoping along the path. Every term on the
second line is nonnegative by \cref{prop:closed-graph}. No bound on the
length of $q$ entered the argument, so the inequality holds for arbitrary
finite input length.
\end{proof}

\begin{remark}
\label{rem:role-of-the-computation}
It is worth separating the two contributions. The mathematics lies in
\cref{lem:state-invariants} and in the telescoping identity
\eqref{eq:exact-endpoint-sum}: the first makes the state set finite, the
second turns three families of inequalities on that finite set into a
statement about words of unbounded length. The computation contributes
only the closed graph and the integer potential, and both are recomputed
from the seeds rather than read from a stored table.
\end{remark}
